\begin{center}
    \bf Resumé
\end{center}

Pour s'adapter aux réseaux internet qui évoluent, le protocole QUIC a été conçu par Google comme un remplacement de TCP, le protocole qui domine aujourd'hui internet.
L'objectif étant de créer un protocole très flexible s'adaptant à toutes les situations, en plus d'être sécurisé et efficace.

Cependant, les implémentations de QUIC étant encore jeunes, leurs performances sont encore loin des implémentations éprouvées de TCP.

Nous nous sommes intéressés à l'une d'elles, Quiche. Après avoir mesuré l'impact d'une version améliorée de Quiche par Florentin Rochet sur cURL, un client HTTP; nous avons réussi à réduire la puissance de calcul consommée par Quiche en ajoutant une nouvelle fonctionnalité.
Nous avons obtenu une réduction significative de l'ordre de 6 à 10\% de la consommation énergétique de Quiche, mais malheureusement la gestion actuelle de certaines erreurs de transmission dans Quiche entre en conflit avec la fonctionnalité que nous avons ajouté et fait chuter le débit d'environ 10\%, ce qui empêche de profiter pleinement de la fonctionnalité que nous avons ajouté.

\vspace*{\stretch{1}}

\begin{center}
    \bf Abstract
\end{center}

To adapt to evolving internet networks, Google designed the QUIC protocol as a replacement for TCP, the protocol that currently dominates the internet.
The goal was to create a highly flexible protocol that could adapt to all situations, in addition to being secure and efficient.

However, as QUIC implementations are still in their infancy, their performance is still far from that of proven TCP implementations.

We took an interest in one of them, Quiche. After measuring the impact of an improved version of Quiche by Florentin Rochet on cURL, an HTTP client, we managed to reduce the computing power consumed by Quiche by adding a new feature.
We achieved a significant reduction of around 6 to 10\% in Quiche's energy consumption, but unfortunately the current handling of certain transmission errors in Quiche conflicts with the feature we added and causes the throughput to drop by around 10\%, which prevents us from taking full advantage of the feature we added.

\vspace*{\stretch{1}}